%=============================================================
% FORM-IC-01 Investigacion Individual
% Arquitecturas Distribuidas en la Industria
% Autor: Vinueza Sanchez Harold Nicolas
% Curso: Aplicaciones Distribuidas - Unidad 3
% Docente: Ing. Guerrero Ulloa Gleiston Ciceron
% Universidad Tecnica Estatal de Quevedo (UTEQ)
% Ciclo academico: 2026-2027 PPA
%=============================================================

\documentclass[11pt,a4paper]{article}

\usepackage[T1]{fontenc}
\usepackage{babel}
\usepackage[margin=2.2cm]{geometry}
\usepackage{parskip}
\usepackage{titlesec}
\usepackage{xcolor}
\usepackage{hyperref}
\usepackage{cite}
\usepackage{fancyhdr}
\usepackage{booktabs}
\usepackage{array}
\usepackage{enumitem}

%--- Color institucional ---
\definecolor{uteqblue}{RGB}{0,70,127}
\definecolor{uteqgray}{RGB}{90,90,90}
\definecolor{rowgray}{RGB}{240,244,248}

%--- Formato de secciones ---
\titleformat{\section}
  {\normalsize\bfseries\color{uteqblue}}
  {\thesection.}{0.5em}{}[\vspace{-2pt}{\color{uteqblue}\rule{\linewidth}{0.6pt}}]
\titlespacing{\section}{0pt}{10pt}{5pt}

%--- Encabezado y pie ---
\pagestyle{fancy}
\fancyhf{}
\renewcommand{\headrulewidth}{0.5pt}
\renewcommand{\headrule}{\hbox to\headwidth{\color{uteqblue}\leaders\hrule height \headrulewidth\hfill}}
\fancyhead[L]{\small\textcolor{uteqblue}{\textbf{UTEQ}} \textcolor{uteqgray}{--- Aplicaciones Distribuidas}}
\fancyhead[R]{\small\textcolor{uteqgray}{FORM-IC-01 $|$ Unidad 3}}
\fancyfoot[C]{\small\textcolor{uteqgray}{\thepage}}

%--- Hiperenlaces ---
\hypersetup{
  colorlinks=true,
  linkcolor=uteqblue,
  citecolor=uteqblue,
  urlcolor=uteqblue,
  pdftitle={FORM-IC-01 - Microservicios en Netflix},
  pdfauthor={Vinueza Sanchez Harold Nicolas}
}

%=============================================================
\begin{document}
%=============================================================

%---------- PORTADA / ENCABEZADO DEL DOCUMENTO ----------
\noindent
\colorbox{uteqblue}{\parbox{\dimexpr\linewidth-2\fboxsep}{%
  \vspace{6pt}
  \centering
  \textcolor{white}{\large\bfseries
    Arquitectura de Microservicios en Producci\'on: Caso Netflix}\\[3pt]
  \textcolor{white!80!uteqblue}{\small
    FORM-IC-01 --- Investigaci\'on Individual $|$ Unidad 3: Arquitecturas Distribuidas en la Industria}
  \vspace{6pt}
}}

\vspace{6pt}
\noindent
\begin{tabular}{@{}ll}
  \textbf{Estudiante:}  & Vinueza S\'anchez Harold Nicol\'as \\
  \textbf{Docente:}     & Ing. Guerrero Ulloa Gleiston Cicer\'on \\
  \textbf{Instituci\'on:} & Universidad T\'ecnica Estatal de Quevedo (UTEQ) \\
  \textbf{Ciclo:}       & Aplicaciones Distribuidas --- 2026--2027 PPA \\
\end{tabular}
\vspace{4pt}
\rule{\linewidth}{0.4pt}

%=============================================================
\section{Descripci\'on del Caso: Netflix y su Transici\'on a Microservicios}

Netflix es hoy la plataforma de \textit{streaming} de video m\'as grande del mundo, con m\'as
de 230 millones de suscriptores en m\'as de 190 pa\'ises y picos de tr\'afico que representan
cerca del 15\,\% del ancho de banda global de Internet. Hasta 2007, toda su infraestructura
resid\'ia en un sistema monol\'itico alojado en centros de datos propios; la totalidad de la
l\'ogica de negocio ---reproducci\'on, recomendaciones, facturaci\'on, autenticaci\'on--- estaba
comprimida en una sola base de c\'odigo que se desplegaba como una unidad indivisible.

En 2008, una corrupci\'on cr\'itica en la base de datos central interrumpi\'o los env\'ios de DVD
durante tres d\'ias consecutivos, exponiendo de forma contundente la fragilidad del modelo
centralizado ante fallos puntuales. Este incidente motiv\'o la toma de decisi\'on m\'as relevante
en la historia de la ingenier\'ia de software de la compa\~n\'ia: migrar progresivamente hacia
Amazon Web Services (AWS) y descomponer el monolito en \textbf{microservicios}
independientes. La transici\'on, que demand\'o aproximadamente siete a\~nos, qued\'o pr\'acticamente
completada hacia 2015~\cite{balalaie2016}.

Hoy, la plataforma opera con m\'as de \textbf{700 microservicios} en producci\'on. Entre los
m\'as cr\'iticos se encuentran el servicio de recomendaci\'on personalizada, el de codificaci\'on y
transcodificaci\'on de v\'ideo (plataforma \textit{Cosmos}), el de autenticaci\'on OAuth, el de
gesti\'on de perfiles y el de facturaci\'on recurrente. Cada uno es propiedad de un equipo
peque\~no (``two-pizza team'') que gestiona su ciclo de vida completo: dise\~no, desarrollo,
pruebas, despliegue y monitoreo~\cite{miguel2022}. Esta estructura organizacional,
conocida como DevOps desentralizado, es una consecuencia directa de la arquitectura
adoptada~\cite{balalaie2016}.

%=============================================================
\section{Por qu\'e Eligieron la Arquitectura de Microservicios}

La decisi\'on de adoptar microservicios en Netflix no fue ideol\'ogica sino pragm\'atica, y
respondi\'o a tres limitaciones estructurales del monolito que se volvieron insostenibles
con el crecimiento exponencial del servicio:

\begin{enumerate}[leftmargin=*, label=\textbf{\arabic{enumi}.}]

  \item \textbf{L\'imite de escalabilidad horizontal.}
  En la arquitectura monol\'itica, escalar un \'unico componente (p.~ej., el motor de
  recomendaciones) implicaba replicar la totalidad del sistema, lo que resultaba en un
  uso ineficiente de recursos y en costos operativos prohibitivos a escala global~\cite{blinowski2022}.

  \item \textbf{Velocidad de entrega comprometida.}
  Un cambio m\'inimo en cualquier m\'odulo obligaba a redesplegar la aplicaci\'on completa,
  con ventanas de mantenimiento superiores a seis horas y tasas de reversi\'on (rollback)
  elevadas. Los microservicios eliminan este acoplamiento al permitir que cada servicio
  tenga su propio pipeline de integraci\'on y despliegue continuo (CI/CD)~\cite{balalaie2016}.

  \item \textbf{Riesgo de fallos en cascada.}
  En un sistema monol\'itico, la falla de un subcomponente no criti\-co pod\'ia derribar la
  experiencia completa del usuario. Los microservicios permiten aislar dominios de negocio
  de modo que los fallos se contienen y no se propagan al sistema global; principio
  formalizado como \textit{bulkhead pattern}~\cite{dragoni2017}.

\end{enumerate}

\noindent
Adicionalmente, Netflix necesitaba una arquitectura que soportara el ciclo de innovaci\'on
continua caracter\'istico de la industria del entretenimiento: nuevas caracter\'isticas
(reproducci\'on sin conexi\'on, co-visualizaci\'on, plan con anuncios) deb\'ian lanzarse en
semanas, no en meses~\cite{miguel2022}.

%=============================================================
\section{Beneficios Obtenidos}

La migraci\'on produjo mejoras sustanciales y cuantificables en tres ejes principales:

\medskip
\noindent\textbf{Escalabilidad independiente y eficiencia de recursos.}
Cada servicio escala en funci\'on exclusiva de su propia demanda. El servicio de
codificaci\'on \textit{Cosmos} puede procesar millones de activos multimedia sin incrementar
los recursos del servicio de autenticaci\'on, que opera con patrones de carga completamente
distintos. Esto redujo significativamente el costo de infraestructura respecto al modelo
monol\'itico~\cite{miguel2022}. Estudios emp\'iricos sobre arquitecturas comparadas confirman
que la escalabilidad granular de los microservicios supera al monolito en escenarios de
carga heterog\'enea~\cite{blinowski2022}.

\medskip
\noindent\textbf{Despliegue continuo y autonomia de equipos.}
La plataforma de Netflix realiza entre 4\,000 y 6\,000 despliegues por d\'ia en producci\'on.
Cada equipo lleva un servicio de \textit{commit} a producci\'on en pocas horas mediante
pipelines automatizados. Esta capacidad es habilitada directamente por el desacoplamiento
arquitect\'onico de los microservicios~\cite{balalaie2016}.

\medskip
\noindent\textbf{Resiliencia demostrada mediante ingenier\'ia del caos.}
Netflix dise\~n\'o el \textit{Chaos Monkey}, una herramienta que inyecta fallos aleatorios en
servicios de producci\'on para validar que el sistema tolera las interrupciones sin
degradar la experiencia del usuario. Esta pr\'actica, hoy extendida a toda la industria
bajo el nombre de \textit{chaos engineering}, fue posible precisamente porque los
microservicios proveen el aislamiento necesario para absorber fallos individuales sin
propagarlos~\cite{dragoni2017}.

%=============================================================
\section{Desaf\'ios Encontrados}

La adopci\'on de microservicios introdujo una clase de complejidad inexistente en el monolito
original. Los principales desaf\'ios documentados en la literatura y confirmados por la
experiencia de Netflix son:

\begin{itemize}[leftmargin=*, itemsep=3pt]

  \item \textbf{Latencia y sobrecarga de red.} Toda interacci\'on entre servicios que
  antes era una llamada en memoria se convierte en una petici\'on de red (HTTP/gRPC), lo
  que introduce latencia medible y requiere estrategias de tolerancia como
  \textit{circuit breakers}, \textit{timeouts} y \textit{retries}~\cite{dragoni2017}.

  \item \textbf{Observabilidad distribuida.} Con m\'as de 700 servicios, identificar la
  causa ra\'iz de un error exige rastreo distribuido (\textit{distributed tracing}),
  correlaci\'on de \textit{logs} entre m\'ultiples instancias y m\'etricas unificadas.
  Netflix desarrollaron herramientas como \textit{Zipkin} y \textit{Atlas} para cubrir
  esta necesidad~\cite{balalaie2016}.

  \item \textbf{Consistencia de datos eventual.} Cada microservicio gestiona su propia
  base de datos, lo que elimina las transacciones ACID globales. Garantizar la
  consistencia entre servicios exige patrones como \textit{sagas} y mecanismos de
  compensaci\'on transaccional~\cite{blinowski2022}.

  \item \textbf{Complejidad operacional.} Orquestar el descubrimiento de servicios,
  balanceo de carga, gesti\'on de secretos y actualizaciones sin tiempo de inactividad
  en cientos de microservicios demanda una plataforma de infraestructura madura
  (Kubernetes, Spinnaker, Consul)~\cite{miguel2022}.

\end{itemize}

%=============================================================
\section{Aplicaci\'on al Proyecto Propio}

Las lecciones del caso Netflix son directamente aplicables al proyecto de la asignatura
Aplicaciones Distribuidas, que contempla el dise\~no de un \textbf{sistema de gesti\'on
acad\'emica distribuido}. A partir del an\'alisis del caso, se identifican las siguientes
decisiones de dise\~no:

\begin{enumerate}[leftmargin=*, itemsep=3pt, label=\textbf{\arabic{enumi}.}]

  \item \textbf{Descomposici\'on por dominio.} Cada funci\'on de negocio ---matr\'icula,
  calificaciones, notificaciones, autenticaci\'on--- se implementa como un microservicio
  independiente con su propia base de datos, siguiendo el principio de
  \textit{bounded context} de Domain-Driven Design~\cite{dragoni2017}.

  \item \textbf{API Gateway centralizado.} Un Gateway \'unico (p.~ej., Kong o AWS API
  Gateway) recibe todas las peticiones externas, gestiona autenticaci\'on, limitaci\'on de
  tasa (\textit{rate limiting}) y enrutamiento, simplificando la interfaz del cliente
  respecto a la complejidad interna del sistema.

  \item \textbf{Comunicaci\'on as\'incrona para operaciones no cr\'iticas.} Procesos como
  el env\'io de notificaciones o la generaci\'on de reportes PDF se desacoplan del flujo
  principal mediante un \textit{message broker} (RabbitMQ o Kafka), mejorando la
  capacidad de respuesta del sistema y su tolerancia a picos de carga~\cite{balalaie2016}.

  \item \textbf{Granularidad controlada.} La lecci\'on m\'as importante del caso Netflix es
  que los microservicios \textit{demasiado finos} generan una red de dependencias que
  anula los beneficios del desacoplamiento. Para el proyecto acad\'emico se propone
  iniciar con 4--5 servicios bien definidos y refinarlos conforme evolucione el
  sistema~\cite{blinowski2022}.

\end{enumerate}

%=============================================================
\vspace{6pt}
{\color{uteqblue}\rule{\linewidth}{0.6pt}}

\begingroup
\small
\begin{thebibliography}{9}

\bibitem{balalaie2016}
A.~Balalaie, A.~Heydarnoori, and P.~Jamshidi,
``Microservices Architecture Enables DevOps: Migration to a Cloud-Native Architecture,''
\textit{IEEE Software}, vol.~33, no.~3, pp.~42--52, May/Jun.\ 2016.
DOI: \href{https://doi.org/10.1109/MS.2016.64}{10.1109/MS.2016.64}

\bibitem{miguel2022}
F.~S.~Miguel, N.~Mareddy, A.~K.~Moorthy, and X.~Liu,
``Microservices for Multimedia: Video Encoding,''
in \textit{Proc.\ 1st ACM Mile-High Video Conf.\ (MHV)}, Denver, CO, USA, 2022.
DOI: \href{https://doi.org/10.1145/3510450.3517298}{10.1145/3510450.3517298}

\bibitem{dragoni2017}
N.~Dragoni, S.~Giallorenzo, A.~L.~Lafuente, M.~Mazzara, F.~Montesi,
R.~Mustafin, and L.~Safina,
``Microservices: Yesterday, Today, and Tomorrow,''
in M.~Mazzara and B.~Meyer (Eds.), \textit{Present and Ulterior Software Engineering},
Springer, Cham, 2017, pp.~195--216.
DOI: \href{https://doi.org/10.1007/978-3-319-67425-4_12}{10.1007/978-3-319-67425-4\_12}

\bibitem{blinowski2022}
G.~Blinowski, A.~Ojdowska, and A.~Przyby{\l}ek,
``Monolithic vs.\ Microservice Architecture: A Performance and Scalability Evaluation,''
\textit{IEEE Access}, vol.~10, pp.~20357--20374, 2022.
DOI: \href{https://doi.org/10.1109/ACCESS.2022.3152803}{10.1109/ACCESS.2022.3152803}

\end{thebibliography}
\endgroup

\end{document}